%% opus_05_G3_holography_topologisation.tex
%%
%% Elite-voice adversarial-heal dossier (Wave 17, 2026-04-24 opus).
%% Target: Identity G3 -- Universal Holography functor
%%   Phi_hol : ChirAlg^{omega, BL}_C  ->  HT-QFT_{C x R}
%% plus the topologisation tower
%%   SC^{ch,top} + inner conformal vector (Sugawara, k != -h^v)  =  E_3-topological
%%   W_N inner stress tower of depth N  ->  E_{N+1}-topological
%%   W_infty  ->  E_infty-topological  (horizon, conditional).
%%
%% Primary voices: Bezrukavnikov (geometric rep theory) + Kazhdan
%% (compression). Secondary: Costello (factorisation).
%%
%% This is a working dossier, not reader-facing prose. Six ATTACK-HEAL
%% cycles; healed identifications; hidden structure; open frontiers;
%% cross-volume harmonies. No manuscript insertion.

\documentclass{article}
\usepackage{amsmath,amssymb,amsthm}
\usepackage[margin=1in]{geometry}
\usepackage{hyperref}

\newtheorem{theorem}{Theorem}[section]
\newtheorem{proposition}[theorem]{Proposition}
\newtheorem{lemma}[theorem]{Lemma}
\newtheorem{corollary}[theorem]{Corollary}
\theoremstyle{definition}
\newtheorem{definition}[theorem]{Definition}
\newtheorem{construction}[theorem]{Construction}
\theoremstyle{remark}
\newtheorem{remark}[theorem]{Remark}
\newtheorem{question}[theorem]{Question}
\newtheorem{attack}[theorem]{Attack}
\newtheorem{heal}[theorem]{Heal}
\newtheorem{hidden}[theorem]{Hidden structure}
\newtheorem{openfrontier}[theorem]{Open frontier}

\newcommand{\SCchtop}{\mathsf{SC}^{\mathrm{ch,top}}}
\newcommand{\Phihol}{\Phi_{\mathrm{hol}}}
\newcommand{\HTQFT}{\mathsf{HT\text{-}QFT}}
\newcommand{\ChirAlg}{\mathsf{ChirAlg}}
\newcommand{\Obs}{\mathsf{Obs}}
\newcommand{\Zder}{Z^{\mathrm{der}}_{\mathrm{ch}}}
\newcommand{\ChirHoch}{\mathrm{ChirHoch}}
\newcommand{\Ehol}{\mathrm{E}_2^{\mathrm{hol}}}
\newcommand{\Etop}[1]{\mathrm{E}_{#1}^{\mathrm{top}}}
\newcommand{\Ech}[1]{\mathrm{E}_{#1}^{\mathrm{ch}}}
\newcommand{\DS}{\mathrm{DS}}
\newcommand{\fg}{\mathfrak{g}}
\newcommand{\Vir}{\mathrm{Vir}}
\newcommand{\cW}{\mathcal{W}}
\newcommand{\cA}{\mathcal{A}}
\newcommand{\cF}{\mathcal{F}}
\newcommand{\MC}{\mathrm{MC}}
\newcommand{\BRST}{\mathrm{BRST}}

\begin{document}
\title{G3 Universal Holography \& Topologisation: six-cycle attack-heal dossier}
\author{Wave 17, opus\_05 (Bezrukavnikov + Kazhdan + Costello voice)}
\date{2026-04-24}
\maketitle

\begin{abstract}
The G3 identity asserts the existence of a canonical functor
$\Phihol\colon \ChirAlg^{\omega,\mathrm{BL}}_C \to
\HTQFT_{C\times\mathbb{R}}$ taking a boundary-linear chiral algebra
with inner conformal vector at non-critical level to a 3D
holomorphic-topological gauge theory on $C\times\mathbb{R}$, with
boundary observables equal to $\cA$ and bulk observables equal to
the derived chiral centre $\Zder(\cA)$. The topologisation theorem
is the engine: the $\SCchtop$ operadic action on the pair
$(C^{\bullet}_{\mathrm{ch}}(\cA,\cA),\cA)$ lifts to
$\Etop{3}$ once an inner conformal vector trivialises the
holomorphic direction on BRST cohomology. Class-L ($V_k(\fg)$,
$k\neq -h^\vee$) is the unconditional base case; class-M
(Virasoro, $\cW_N$) is weight-completed only; the original-complex
chain-level statement for class-M is the genuine open frontier.
This dossier walks through six ATTACK-HEAL cycles, constructs
$\Phihol$ on $V_1(\mathfrak{sl}_2)$ from first principles, extends
to $\cW_3$ as the $N=3$ ladder rung, and maps the $\cW_\infty$
horizon to three independently checkable convergence criteria.
\end{abstract}

\tableofcontents

%=======================================================================
\section{Setup and the seven attack questions}
\label{sec:opus-setup}
%=======================================================================

The target identity is G3 in the Platonic reconstitution
(file `platonic_ideal_reconstituted_2026_04_17.md`, \S3):
\begin{equation}\label{eq:opus-G3-statement}
  \Phihol\colon \ChirAlg^{\omega,\mathrm{BL}}_C
  \longrightarrow \HTQFT_{C\times\mathbb{R}},
  \qquad
  (\cA,\omega,\mathrm{BL}) \longmapsto
  (\mathrm{boundary}=\cA,\;\mathrm{bulk}=\Zder(\cA),\;
   \mathrm{interaction}=\SCchtop\text{-braces}),
\end{equation}
together with the topologisation theorem
\begin{equation}\label{eq:opus-topologisation}
  \SCchtop + \omega\ \text{(inner, non-critical)}
  \;\longrightarrow\; \Etop{3}
  \quad\text{on}\quad \Zder(\cA).
\end{equation}
The inscribed source files are:
\begin{itemize}
\item Vol II master theorem: \texttt{chapters/connections/%
universal\_holography\_functor.tex},
Thm.~\ref*{opus:thm:UHF}.
\item Vol II pointwise master + $c=24$ landscape plus Banach radius:
\texttt{chapters/connections/programme\_climax\_platonic.tex},
Thm.~\ref*{opus:thm:PMC}.
\item Vol II $\Einfty$-ladder: \texttt{chapters/connections/%
e\_infinity\_topologization.tex}, Thm.~\ref*{opus:thm:EINF-LAD}.
\item Vol I topologisation base: \texttt{chapters/theory/%
en\_koszul\_duality.tex}, \S\ref*{opus:sec:V1topo}.
\item Vol II class-M chain-level status:
\texttt{chapters/theory/topologization\_class\_m\_%
original\_complex\_platonic.tex}.
\end{itemize}

\begin{question}[Seven attack questions]
\label{q:opus-seven}
\begin{enumerate}[label=\textup{(\alph*)}]
\item \textbf{Identification vs lift.} Is the topologisation theorem
an \emph{identification} of $\SCchtop + \omega$ with $\Etop{3}$, or
only a \emph{lift} up to weight-completion?
\item \textbf{``Inner'' conformal vector.} Is ``inner'' a structural
requirement or just any $V_\omega \in \cA$ of weight two?
\item \textbf{Regularity of the Sugawara topologisation.} At what
regularity does the Sugawara construction topologise the
$C$-direction? Does the argument fail at $k=-h^\vee$?
\item \textbf{$\cW_N$ chain-level status.} Is the
$\cW_N \to \Etop{N+1}$ statement proved chain-level, or only at
weight-completion?
\item \textbf{$\cW_\infty$ horizon.} Is the
$\cW_\infty \to \Etop{\infty}$ claim a genuine operadic horizon, or
does the tower terminate finitely?
\item \textbf{Scope of $\Phihol$.} Is
``non-logarithmic $C_2$-cofinite'' the maximal unconditional scope?
\item \textbf{Commit \texttt{a5640de}.} Did the \texttt{a5640de}
commit close $H^2=0$ at all genus, or only at specific genus?
\end{enumerate}
\end{question}

These seven questions drive six ATTACK-HEAL cycles
(\S\ref{sec:opus-cycle-1} through \S\ref{sec:opus-cycle-6}).
Attack~(g) is dismissed inside Cycle~1 since a direct
\texttt{git show a5640de} reveals the commit is a $W(p)$-tempered
retraction, unrelated to $H^2$. Attacks~(a)--(f) carry the
mathematics.

%=======================================================================
\section{Cycle 1: identification versus lift; the weight-completion axis}
\label{sec:opus-cycle-1}
%=======================================================================

\begin{attack}[a,~g]
\label{att:cycle1}
The programme sometimes writes
``$\SCchtop + \omega = \Etop{3}$''
(master-theorem statements in
\texttt{universal\_holography\_functor.tex} and
\texttt{programme\_climax\_platonic.tex}). Two objections:
\begin{enumerate}[label=\textup{(\roman*)}]
\item The equality sign hides an $(\infty,1)$-categorical
promotion. $\SCchtop$ is a two-coloured operad (or
\emph{bicoloured propad}; see
\texttt{class\_m\_direct\_sum\_obstruction\_platonic.tex}); the
promotion to $\Etop{3}$ is the content of Dunn additivity
\emph{after} the holomorphic direction has been topologised by
an inner conformal vector. Dunn is a theorem about
$\infty$-operads in the topological monoidal
$\infty$-category; calling it an ``equality'' conflates
chain-level data with $(\infty,1)$-categorical identification.
\item For class-M the inscribed chain-level statement only
holds on the \emph{weight-completed} complex
(\texttt{prop:bv-bar-class-m-weight-completed});
the original-complex chain-level statement is a frontier item.
Commit \texttt{a5640de} did not close this. (\texttt{git show
a5640de} shows the commit is the $W(p)$-tempered retraction,
unrelated to $H^2$ or the class-M original-complex statement.)
\end{enumerate}
\end{attack}

\begin{heal}[a,~g]
\label{heal:cycle1}
Both objections are absorbed into a three-tier layer structure
(Vol~I~\S\ref*{opus:sec:V1topo},
Thm.~\ref*{opus:thm:topologization}):

\smallskip
\textbf{Tier~I (cohomological identification).} On BRST
cohomology $H^\bullet_Q\bigl(\Zder(\cA)\bigr)$, once $\omega$
acts $Q$-exactly on holomorphic translations
($T_{\mathrm{Sug}} = [Q, G]$), the factorisation algebra on the
holomorphic direction becomes locally constant, which by
Lurie HA Theorem~5.4.5.9 recognises an $\Etop{2}$ structure.
Dunn additivity with the transverse $\Etop{1}$ then gives
\emph{genuine equality of operads}:
\begin{equation}\label{eq:opus-tier-I}
  H^\bullet_Q\bigl(\Zder(\cA)\bigr)
  \;\in\; \mathrm{Alg}_{\,\Etop{2} \otimes_{\mathrm{Dunn}} \Etop{1}}
  \;=\; \mathrm{Alg}_{\,\Etop{3}}.
\end{equation}
This is an identification, not a lift. Tier~I is
unconditional for class~L (affine KM) and
class-M-at-cohomology, at every genus.

\smallskip
\textbf{Tier~II (chain-level, transferred model).} Choose a
strong deformation retract
$(\pi, \iota, h)\colon \Zder(\cA) \rightleftarrows
H^\bullet_Q\bigl(\Zder(\cA)\bigr)$ (Cauchy datum,
\texttt{homotopy retract is data}, RS-2). Transport the
Tier~I structure across the retract via the homotopy transfer
theorem (Kontsevich--Soibelman, Lada--Markl), producing a
chain-level $\Etop{3}_\infty$-algebra structure on the
zero-differential model. The result is unique up to
$\infty$-quasi-isomorphism. Tier~II is unconditional for the
same scope as Tier~I.

\smallskip
\textbf{Tier~III (original-complex chain-level).} The
transferred structure of Tier~II is on the cohomology model,
not on $\Zder(\cA)$ itself. To obtain an $\Etop{3}$
structure on the \emph{original} chain complex requires an
$A_\infty$-coherent null-homotopy of holomorphic translation
in the brace deformation complex
(\texttt{en\_koszul\_duality.tex} \S\ref*{opus:sec:V1topo},
Part~(iii) of Thm.~\ref*{opus:thm:topologization}): an
$A_\infty$-family $G = (G_1, G_2, \ldots)$ with
$[m, G] = \partial_z$, $G_1 =$ Sugawara antighost
contraction. For class~L this has been closed
(Prop.~\texttt{prop:sugawara-gauge-rectification}); the
tower terminates at $r_{\max} = 3$. For class-M the
tower need not terminate, and only the weight-completed
statement is proved.

\smallskip
\textbf{Scope declaration.} The G3 statement
``$\Phihol$ exists as a functor'' holds in the
$(\infty,1)$-categorical lane at Tier~I (cohomological). The
chain-level lane (Tiers~II, III) refines the same statement
at chain level. Tier~II is proved unconditionally across
G/L/C and class-M-weight-completed; Tier~III is proved
unconditionally for class~L and is a frontier item for
class~M original complex.

\smallskip
\textbf{Commit \texttt{a5640de}.} \texttt{git show a5640de}
(2026-04-17, author Raeez Lorgat) retracts
\texttt{thm:tempered-stratum-contains-wp} for the logarithmic
triplet $W(p)$: it downgrades the status from
$\textsf{ProvedHere}$ to $\textsf{Conjectured}$ after finding
that the Zhu-bounded-Massey argument fails on primary-source
audit (Gurarie 1993, Flohr 1996). The commit has no
relationship to $H^2$, curved-Dunn, or class-M
original-complex. The G3 attack question (g) is dismissed.
\end{heal}

\begin{hidden}[layer structure]
\label{hidden:cycle1}
The three tiers are not a hierarchy in which higher tiers
subsume lower ones; they are three different theorems about
three different mathematical objects: BRST cohomology, a
cohomology-model quasi-iso, and the original chain complex.
Each tier is loadbearing, and the G3 identity in its
\emph{strongest} honest form is the statement that
$\Phihol$ lands on $\Etop{3}$-topological algebras at
Tier~I globally, at Tier~II via transferred models
globally, and at Tier~III on the restricted subcategory of
class-L + class-M-weight-completed objects.
\end{hidden}

%=======================================================================
\section{Cycle 2: what ``inner'' means (structural requirement)}
\label{sec:opus-cycle-2}
%=======================================================================

\begin{attack}[b]
\label{att:cycle2}
The hypothesis in the source category
$\ChirAlg^{\omega,\mathrm{BL}}_C$
(\texttt{universal\_holography\_functor.tex}
Def.~\ref*{opus:def:UHF-source}) is that $\omega \in \cA_2(X)$
is a conformal vector and $\cA$ is boundary-linear. There is
an additional hypothesis that $\omega$ is an
\emph{inner} conformal vector at a non-critical level
(Definition~\ref{def:opus-inner}). Is ``inner''
\emph{needed}, or is any weight-2 Virasoro-OPE generator
enough? If merely Virasoro-OPE is required, the source
category is much larger than stated; if ``inner''' is
structural, the statement needs this clause to close.
\end{attack}

\begin{heal}[b]
\label{heal:cycle2}
``Inner'' is structural, not decorative. Here is why.

\begin{definition}[Inner conformal vector]
\label{def:opus-inner}
Let $\cA$ be a chiral algebra with a weight-2 generator
$\omega$ satisfying the Virasoro OPE of central charge~$c$.
Let $Q$ be the BRST differential of a 3d
holomorphic-topological theory $T_\cA$ on $C \times \mathbb{R}$
whose boundary chiral algebra is $\cA$. The conformal vector
$\omega$ is \emph{inner} for $(\cA, Q)$ if there exists
$G(z)$ in the bulk observables of $T_\cA$ such that
\begin{equation}\label{eq:opus-inner}
  \omega(z) \;=\; [Q, G(z)]
  \quad\text{in }\ H^\bullet\bigl(A^{\mathrm{BV}}_{3d},Q\bigr).
\end{equation}
\end{definition}

\begin{proposition}[Failure without ``inner'';
\texttt{proved elsewhere}: \texttt{en\_koszul\_duality.tex}
Rem.~\ref*{opus:rem:conformal-vector-necessary}]
\label{prop:opus-need-inner}
If $\omega$ is not inner for the bulk $Q$, then the
topologisation step fails at Tier~I: $\partial_z$ is not
$Q$-exact on $H^\bullet_Q\bigl(\Zder(\cA)\bigr)$, and the
factorisation algebra on the holomorphic direction retains
genuine $\bar\partial$-dependence. The $\Etop{2}$ structure
on $\Zder(\cA)$ produced by chiral Deligne stays
\emph{holomorphic} $\Ehol$, not topological, and the
Dunn-additivity step $\Etop{2} \otimes \Etop{1} \simeq
\Etop{3}$ does not apply.
\end{proposition}

\begin{proof}[Sketch]
Without innerness there is no BRST-primitive $G$ lifting
$\omega$. The Sugawara identity
$\partial_z = [T, \cdot] = [[Q, G], \cdot] = [Q, [G, \cdot]]$
breaks at the second step. The factorisation algebra on
the $C$-direction therefore carries $\bar\partial$-cohomology
classes at every point of $C$, so local constancy fails.
\end{proof}

\begin{hidden}[inner $\leftrightarrow$ bulk BV-BRST existence]
\label{hidden:cycle2}
Innerness is simultaneously a \emph{property of $\cA$} (the
existence of the BRST primitive $G$) and a \emph{construction
of the bulk theory} (the existence of the
$A^{\mathrm{BV}}_{3d}$ complex). The two are inseparable: one
cannot write down $G$ without specifying the bulk theory, and
one cannot construct the bulk theory without a chiral
boundary algebra $\cA$ that is the Dirichlet / boundary
condition for the BV fields. For the four standard
landscape families:
\begin{itemize}
\item \textbf{G} (Heisenberg / lattice): $\omega =
\tfrac{1}{2}\sum \alpha \alpha$ is inner, with
$G = \tfrac{1}{2}\sum \bar c_\alpha \alpha$ the abelian
antighost bilinear. Bulk: Costello--Li abelian holomorphic
Chern--Simons.
\item \textbf{L} (affine KM $V_k(\fg)$, $k \neq -h^\vee$):
$\omega = T_{\mathrm{Sug}}$ is inner, with
$G = \tfrac{1}{2(k+h^\vee)} \sum :J^a \bar c_a:$. Bulk:
Costello--Gaiotto non-abelian holomorphic Chern--Simons.
\item \textbf{C} ($\beta\gamma$, $bc$): bosonised to
class~G; inherits innerness.
\item \textbf{M} (Virasoro, $\cW_N$, hook-type
$\cW_k(\fg, f)$): $\omega = T$ (stress tensor). Inner for
the Costello--Gaiotto bulk with Drinfeld--Sokolov boundary;
the DS-reduction provides the bulk
$A^{\mathrm{BV}}_{3d}$-complex and the primitive $G$.
\end{itemize}
The B-row (BKM extension with Borcherds product; Mukai-enhanced
K3 Heisenberg) has a \emph{non-inner} conformal vector
(the Borcherds weight $5$ does not match $T_{\mathrm{Sug}} + $
Heisenberg decomposition): the topologisation route on the
B-row goes through the direct Mukai-enhanced K3 Heisenberg
presentation (Vol~III, \texttt{cy\_d\_kappa\_stratification%
.tex}) rather than through Sugawara.
\end{hidden}

%=======================================================================
\section{Cycle 3: Sugawara regularity and critical-level failure}
\label{sec:opus-cycle-3}
%=======================================================================

\begin{attack}[c]
\label{att:cycle3}
What is the \emph{regularity} of the Sugawara topologisation
on the $C$-direction? Is it chain-level, or only after
completion? The Sugawara polynomial
$T_{\mathrm{Sug}} = \frac{1}{2(k+h^\vee)}\sum_a {:}J^a J_a{:}$
involves a level-dependent denominator that blows up at
$k = -h^\vee$. Does the argument fail at the critical level,
and is this a genuine restriction or a convention issue?
\end{attack}

\begin{heal}[c]
\label{heal:cycle3}
The regularity is \emph{single-contour normal-ordered chain-
level} for all $k \neq -h^\vee$ and \emph{fails genuinely} at
$k = -h^\vee$. The healed picture:

\textbf{Non-critical regularity.} For any non-critical level
$k \neq -h^\vee$, the Sugawara polynomial
\begin{equation}\label{eq:opus-sugawara}
  T_{\mathrm{Sug}}(z) \;=\;
  \frac{1}{2(k+h^\vee)}\, \sum_a {:}J^a(z) J_a(z){:}
\end{equation}
is a well-defined single-contour normally ordered
field in $\cA$ (Kac \S5; Frenkel--Ben-Zvi \S4.5). Its modes
$L_n = \mathrm{Res}_z z^{n+1} T_{\mathrm{Sug}}(z)$ satisfy
the Virasoro OPE of central charge
\begin{equation}\label{eq:opus-ksug-c}
  c_{\mathrm{Sug}}(k) \;=\;
  \frac{k \cdot \dim(\fg)}{k+h^\vee}.
\end{equation}
The BRST primitive $G$ of Cycle~2 exists chain-level in the
3d BV complex, and the identity $T_{\mathrm{Sug}} = [Q, G]$
holds strictly on BRST cohomology (Tier~I); for class~L,
also strictly on the original chain complex (Tier~III, via
gauge rectification).

\textbf{Critical-level failure.} At $k = -h^\vee$:
\begin{enumerate}[label=(\arabic*)]
\item The denominator in \eqref{eq:opus-sugawara} vanishes,
so $T_{\mathrm{Sug}}$ is undefined as a normally ordered
polynomial in $J$-currents;
\item The Feigin--Frenkel centre
$\mathfrak{z}(\hat\fg)_{-h^\vee}$ of the universal affine
algebra becomes commutative (not Sugawara!);
the Sugawara mechanism for topologisation collapses;
\item The correct replacement at critical level is the
Feigin--Frenkel isomorphism
$\mathfrak{z}(\hat\fg)_{-h^\vee} \simeq
\mathrm{Fun}(\mathrm{Op}_{\fg^L}(D))$ identifying the centre
with the space of opers on the formal disc for the Langlands
dual; but this is \emph{arithmetic} Koszul duality
(Beilinson--Drinfeld), not Sugawara topologisation. The bulk
theory becomes wild (irregular-singular KZB, Stokes data);
the $\Etop{3}$ upgrade fails.
\end{enumerate}

\begin{remark}[Not a convention issue, a genuine restriction]
\label{rem:opus-critical-scope}
The hypothesis $k \neq -h^\vee$ in the source category
$\ChirAlg^{\omega,\mathrm{BL}}_C$ is structural, not a
convention one can rephrase away. The master theorem of
$\Phihol$ is \emph{undefined} on
$V_{-h^\vee}(\fg)$. This is consistent with the general
principle that critical-level affine algebras belong to a
different geometric setting (geometric Langlands, opers,
Drinfeld--Sokolov BRST) than non-critical-level ones
(holomorphic Chern--Simons, $\Etop{3}$-topological).
\end{remark}

\begin{hidden}[critical $\leftrightarrow$ arithmetic duality]
\label{hidden:cycle3}
What $\Phihol$ assigns at non-critical level is a
\emph{3D quantum gauge theory}. At critical level the natural
assignment is the \emph{oper-moduli} of the Langlands-dual
group. The two assignments live on different sides of the
Langlands wall: one on the ``automorphic side'' (geometric
Langlands with Kac--Moody symmetry, opers, Hecke modifications);
the other on the ``spectral side''
($\Etop{3}$-topological local observables, quantum gauge
theory). $\Phihol$ is the functor on the spectral side only;
the critical-level analogue is a separate theorem
(Beilinson--Drinfeld's geometric Langlands programme).
\end{hidden}
\end{heal}

%=======================================================================
\section{Cycle 4: $\cW_N$ chain-level and the $N=3$ test case}
\label{sec:opus-cycle-4}
%=======================================================================

\begin{attack}[d]
\label{att:cycle4}
The $\Einfty$-ladder theorem
(\texttt{e\_infinity\_topologization.tex}
Thm.~\ref*{opus:thm:EINF-LAD}) claims
$\cW_N \to \Etop{N+1}$
for every depth~$N$, with $\cW_\infty \to \Etop{\infty}$ in
the limit. Two worries:
\begin{enumerate}[label=\textup{(\roman*)}]
\item Is this a \emph{chain-level} statement on the original
complex, or a weight-completed one? For class~M the
general pattern (Cycle~1) says only weight-completed;
does that apply to $\cW_N$?
\item What is the scope of the
\texttt{tower:antighost-commutativity} axiom (axiom T5 of
\texttt{e\_infinity\_topologization.tex})? The inscribed
remark \texttt{rem:axiom-T5-scope} says it is vacuous at
$N=2$ and reducible to a single bracket
$[G^{(2)}, G^{(3)}]$ at $N=3$, but \emph{postulated} for
$N \ge 4$. Where exactly does the proof become conditional?
\end{enumerate}
\end{attack}

\begin{heal}[d]
\label{heal:cycle4}
The healed status is precise. We give the step-by-step
first-principles construction of $\Phihol$ on
$\cW_3^k$ (the Zamolodchikov algebra, $\fg = \mathfrak{sl}_3$,
principal nilpotent) through the $N=3$ ladder rung.

\textbf{Stress tower of $\cW_3$.} The Zamolodchikov
algebra at generic $c$ has two primary generators: the
stress tensor $T(z)$ of conformal weight $2$ and the
Zamolodchikov field $W(z)$ of conformal weight $3$
(Zamolodchikov 1985; Fateev--Lukyanov 1988). Axioms
T1--T4 of the higher-spin stress tower
(\texttt{def:higher-spin-stress-tower}) are verified:
\begin{itemize}
\item T1 (quasiprimary): both $T$ and $W$ are
quasiprimary w.r.t.\ the Virasoro subalgebra generated by
$T$.
\item T2 (inner Sugawara-like): $T = T_{\mathrm{Sug}}
+ T_{\mathrm{imp}}$, with
$T_{\mathrm{imp}} = (\rho, \partial H)$ the principal
improvement; $W$ is the cubic Casimir
$W = d_{abc}^{(3)} :J^a J^b J^c: + \text{improvements}$,
the Fateev--Lukyanov explicit formula.
\item T3 (truncated $\cW_\infty$): the OPE of $\{T, W\}$
closes on $\{T, W, :WW:, :TW:, \partial T, \partial W, \ldots\}$,
a quotient of Linshaw's universal $\cW_\infty[\mu]$ at
truncation depth 3.
\item T4 (BRST-primitive existence): the 3D BV complex for
$\cW_3^k$ is the Costello--Gaiotto bulk with principal DS
boundary on $\widehat{\mathfrak{sl}}_3$. The
antighost $G^{(2)} = \tfrac{1}{2(k+3)} \sum_a :J^a \bar c_a:$
gives $T = [Q_{\mathrm{tot}}, G^{(2)}]$ on cohomology
(Cycle~2). For $G^{(3)}$, the higher-spin
Casimir construction
(\texttt{eq:G-n-formula}) gives
\begin{equation}\label{eq:opus-G3}
  G^{(3)}(z) \;=\; \frac{1}{3} \sum_{j=1}^3
  d^{a_1 a_2 a_3}\, :J_{a_1}(z)\cdots \bar c_{a_j}(z)
  \cdots J_{a_3}(z):
\end{equation}
with $d$ the totally symmetric invariant tensor of rank 3 on
$\mathfrak{sl}_3$; the identity $W = [Q_{\mathrm{tot}},
G^{(3)}]$ on cohomology follows by Leibniz
(\texttt{eq:leibniz-tower}) and
$J_a = [Q_{\mathrm{tot}}, \bar c_a]$
(Costello--Gaiotto boundary condition).
\end{itemize}

\textbf{Axiom T5 at $N=3$ (BRST-commutativity of
antighosts).} The single bracket to verify is
$[G^{(2)}, G^{(3)}]$. By Leibniz:
\begin{align*}
  [G^{(2)}, G^{(3)}]
  &= \left[\tfrac{1}{2(k+3)}\sum :J^a \bar c_a:,\ G^{(3)}\right] \\
  &= \tfrac{1}{2(k+3)} \sum \left(\text{terms with one
$J$ replaced by $[:J\bar c:, \cdot]$-contraction}\right).
\end{align*}
Each contraction reduces to a combination of single-contour
OPE residues between $J^a$ and the $J$-insertions inside
$G^{(3)}$. By the classical Poisson identity (Linshaw
Thm.~7.1; BMP \S4.3), $\{T^{(2)}, T^{(3)}\} =$ total
derivative of a spin-$2$ combination of $T^{(2)}$ and
$T^{(3)}$; lifting to the quantum OPE via
Feigin--Frenkel screening gives
$[T^{(2)}, T^{(3)}] = \partial f$ on cohomology for some
quantum field $f$. Substituting $T^{(n)} = [Q,G^{(n)}]$ gives
$[G^{(2)}, G^{(3)}] = Q$-exact. Axiom T5 verified at $N=3$.

\textbf{Ladder at $N=3$.} Apply
Thm.~\ref*{opus:thm:EINF-LAD}:
\begin{equation}\label{eq:opus-W3-ladder}
  \Zder(\cW_3^k) \;\in\;
  \mathrm{Alg}_{\,\Etop{2}(G^{(2)})
    \otimes \Etop{1}(G^{(3)})
    \otimes \Etop{1}_{\mathrm{tr}}}
  \;=\; \mathrm{Alg}_{\,\Etop{4}}.
\end{equation}
Dimension count: $2$ (from $G^{(2)}$ topologising the
$\mathbb{C}$-direction) $+\ 1$ (from $G^{(3)}$ topologising
the first-order jet) $+\ 1$ (transverse $\mathbb{R}$) $= 4$.

\textbf{Chain-level scope.} Tier~I (cohomological) is
unconditional. Tier~II (transferred model) is unconditional.
Tier~III (original complex): the analogue of the class~L
gauge rectification is \emph{open} for $\cW_3$ on the
uncompleted direct-sum complex; the weight-completed
version holds by the DS-Hochschild bridge to
$V_k(\mathfrak{sl}_3)$. This is the class-M weight-completed
scope of Cycle~1 instantiated at $\cW_3$.

\textbf{General $\cW_N$.} The construction generalises: the
$N-1$ Casimir tensors $T^{(2)}, T^{(3)}, \ldots, T^{(N)}$
each admit antighost primitives $G^{(n)}$ by
\eqref{eq:opus-G3}-type formulas. Axiom T5 at $N$ requires
$(N-1)(N-2)/2$ pairwise brackets
$[G^{(m)}, G^{(n)}]$ for $2 \le m < n \le N$ to be
$Q$-exact. At $N=3$: one bracket, verified. At $N=4$: three
brackets, \emph{postulated}. The classical Poisson shadow of
these brackets \emph{does} vanish modulo total derivatives
of lower-spin Casimirs (Linshaw Thm.~7.1); the lift from
classical Poisson to quantum BRST-commutation is what
remains to verify from first principles at $N \ge 4$.

\begin{hidden}[$N=3$ as the witness where all inputs line up]
\label{hidden:cycle4}
The $\cW_3$ case is the minimal non-trivial verification of
the ladder beyond the base rung $N=2$. At $N=3$: the
bulk $3$d theory is Costello--Gaiotto holomorphic
Chern--Simons on $\widehat{\mathfrak{sl}}_3$ with principal
DS boundary; the stress tower has the classically-computable
cubic Casimir $W$ with Fateev--Lukyanov formula; axiom T5
reduces to one testable OPE bracket; the resulting bulk
algebra is $\Etop{4}$, matching the Zamolodchikov conformal
field theory's interpretation as 4D higher-spin gravity in
the Gaberdiel--Gopakumar dictionary.
\end{hidden}
\end{heal}

%=======================================================================
\section{Cycle 5: the $\cW_\infty$ horizon}
\label{sec:opus-cycle-5}
%=======================================================================

\begin{attack}[e]
\label{att:cycle5}
Is $\cW_\infty \to \Etop{\infty}$ a genuine operadic horizon
or does the tower terminate at some finite depth? The naive
reading of the formula
``$\Etop{N+1} \otimes_{\mathrm{Dunn}} \Etop{1} =
\Etop{N+2}$'' suggests that every Sugawara step genuinely
increases the dimension, but each truncation
$\cW_\infty \twoheadrightarrow \cW_N$ drops infinitely many
higher-spin generators. Why doesn't the ladder stabilise?
\end{attack}

\begin{heal}[e]
\label{heal:cycle5}
The tower does not terminate because each higher-spin
Casimir is \emph{independent} (neither a polynomial in
lower-spin generators nor a total derivative). The
$\Einfty$-ladder is a genuine operadic horizon, realised
as an inverse limit in the $\infty$-operad category.

\textbf{Operadic statement.} By
Lurie HA Notation~5.1.1.5, $\Etop{\infty}$ is \emph{defined}
as the inverse limit of the stabilisation tower
\begin{equation}\label{eq:opus-Einfty-def}
  \Etop{\infty} \;=\; \varprojlim_n \Etop{n}
\end{equation}
in the symmetric monoidal $\infty$-category of
$\infty$-operads
$\mathrm{Op}_\infty^{\otimes}$. The $\cW_\infty$ horizon
statement is
\begin{equation}\label{eq:opus-Winfty-horizon}
  \Zder(\cW_\infty[\mu]) \;\in\;
  \mathrm{Alg}_{\,\Etop{\infty}}
  \;=\; \varprojlim_N \mathrm{Alg}_{\,\Etop{N+1}},
\end{equation}
where the inverse limit is along the truncation tower
$\cW_\infty[\mu] \twoheadrightarrow \cW_N[\mu]$ of Linshaw.

\textbf{Independence of higher Casimirs.} At generic $\mu$,
the algebra $\cW_\infty[\mu]$ contains a quasiprimary field
$T^{(n)}$ of every spin $n \ge 2$, and these are
algebraically independent from $\{T^{(2)}, \ldots,
T^{(n-1)}\}$: specifically, Linshaw shows that the commutant
$\mathrm{Com}(\cW_{N-1}[\mu], \cW_\infty[\mu])$ is
generated by $T^{(n)}$ for $n \ge N$, strictly strictly
containing fields for every $n$. Therefore each step in the
ladder contributes a genuinely independent $\Etop{1}$-factor,
and the Dunn product does not stabilise.

\textbf{Three equivalent convergence criteria.}
(\texttt{conj:w-infty-e-infty-topological-convergence},
\texttt{e\_infinity\_topologization.tex}):
\begin{enumerate}[label=\textup{(\alph*)}]
\item \textbf{Gaberdiel--Gopakumar stabilisation.} Every OPE
structure constant $f^{n_1 n_2}_{n_3, k}(c)$ of
$\cW_\infty[\mu]$ equals the corresponding structure constant
of $\cW_N[\mu]$ for $N \ge \max(n_1, n_2, n_3, k)$.
\item \textbf{$\infty$-categorical Dunn convergence.} The
iterated Dunn product $\bigotimes_{n \ge 2} \Etop{1}(G^{(n)})$
converges in $\mathrm{Op}_\infty^{\otimes}$ to
$\Etop{\infty}$.
\item \textbf{Yamada strong convergence.} On every
bar-weight window $[w_{\min}, w_{\max}]$, the truncation
$\cW_\infty[\mu] \twoheadrightarrow \cW_N[\mu]$ induces
isomorphisms of weight-$w$ subspaces for $N \ge w_{\max} - 1$.
\end{enumerate}

Criterion (a) is proved by Gaberdiel--Gopakumar
(\texttt{GaberdielGopakumar2011-Winfty} Thm.~3.1).
Criterion (c) is proved by Yamada
(\texttt{Yamada-Winfty-convergence}) at generic $c$.
Criterion (b) follows from (c) by co-cofinality of the
stabilisation tower (Ayala--Francis Lemma~3.7).

\textbf{Scope of the $\cW_\infty$ horizon statement.}
Conditional on axiom T5 at all spins (Cycle~4), the
$\Einfty$-structure on
$\Zder(\cW_\infty[\mu])$ is realised chain-level via the
weight-filtered inverse limit. Each finite weight window
$[w_{\min}, w_{\max}]$ is computed in a finite truncation
$\cW_{w_{\max}-1}[\mu]$ where the $\Etop{w_{\max}}$
structure is Tier~II-transferred. The horizon statement is
therefore:
\begin{equation}\label{eq:opus-Winfty-conditional}
  \text{Tier~I: unconditional at generic }c;\quad
  \text{Tier~II: conditional on T5 at all spins};\quad
  \text{Tier~III: OPEN.}
\end{equation}

\begin{hidden}[horizon $=$ Koszul self-dual endpoint]
\label{hidden:cycle5}
$\Etop{\infty}$ is Koszul self-dual at the operadic level
(Fresse, Thm.~14.1.1): the Koszul dual of $\Etop{\infty}$
is itself. The $\cW_\infty[\mu]$ endpoint is therefore the
\emph{unique} standard-landscape chiral algebra whose
derived chiral centre lies on the self-dual diagonal of the
operadic E-ladder. This matches the architectural
expectation that $\cW_\infty[\mu]$ is the Platonic endpoint
of all higher-spin symmetry: it is what the
Gaberdiel--Gopakumar AdS/CFT correspondence calls the
``higher-spin square'' boundary algebra, whose bulk is
Vasiliev's higher-spin gravity in $\mathrm{AdS}_3$.
\end{hidden}
\end{heal}

%=======================================================================
\section{Cycle 6: scope of unconditional $\Phihol$}
\label{sec:opus-cycle-6}
%=======================================================================

\begin{attack}[f]
\label{att:cycle6}
The master theorem
(\texttt{programme\_climax\_platonic.tex}
Thm.~\ref*{opus:thm:PMC}) asserts $\Phihol$ unconditionally
on the ``non-logarithmic, $C_2$-cofinite standard
landscape'' plus ``VSKR+BGG-tempered irrational cosets''.
Is this the maximal scope? Two directions of worry:
\begin{enumerate}[label=\textup{(\roman*)}]
\item Why ``non-logarithmic''? What fails for
logarithmic $W(p)$?
\item Why ``$C_2$-cofinite''? What fails for irrational
non-tempered cosets?
\end{enumerate}
\end{attack}

\begin{heal}[f]
\label{heal:cycle6}
\textbf{Why ``non-logarithmic''.} Logarithmic vertex
algebras (the triplet $W(p)$ being the canonical example)
have OPEs that carry $\log(z-w)$ singularities. Four
structural failures ensue:
\begin{itemize}
\item $L_0$ is non-diagonalisable: it has a Jordan block of
rank 2 on the top-level fields (the triplet and its
logarithmic partner). This breaks the
``quasiprimary'' axiom T1 of the higher-spin stress tower.
\item The Zhu algebra of $W(p)$ is not semisimple; the
Zhu-bounded-Massey argument of commit \texttt{a5640de}
fails precisely here (Gurarie 1993, Flohr 1996:
logarithmic Massey products grow super-polynomially despite
$C_2$-cofiniteness).
\item The tempered stratum of the $\Phihol$-image does not
contain $W(p)$, as proved by the retraction of
\texttt{thm:tempered-stratum-contains-wp}.
\item The bulk theory is expected to be logarithmic 3D
gravity (Grumiller--Riedler 2010), whose BRST complex has
an $L_0$-nilpotent sector not captured by the single
Sugawara mechanism.
\end{itemize}

\textbf{Why ``$C_2$-cofinite''.} $C_2$-cofiniteness is
Zhu's finiteness condition: the quotient
$\cA / C_2(\cA)$ by the second derivation filtration is
finite-dimensional. Three structural uses:
\begin{itemize}
\item $C_2$-cofinite implies finite-dim Zhu algebra
$\mathrm{Zhu}(\cA)$ (Zhu 1996 Thm.~4.4.3);
\item $C_2$-cofinite + rational implies modular invariance
of characters (Zhu);
\item $C_2$-cofinite implies the chiral algebra has
\emph{finite shadow depth} (Arakawa): the shadow
obstruction tower has $r_{\max} < \infty$. This is the
scope on which $\Phihol$ produces an $\Etop{3}$-topological
rather than a more exotic bulk.
\end{itemize}

\textbf{Irrational cosets.} Irrational cosets
$\mathrm{Com}(B, C)$ need not be $C_2$-cofinite but are
still \emph{tempered} in the VSKR+BGG sense
(\texttt{tempered\_stratum\_characterization\_platonic.tex}):
their shadow tower, while infinite, admits a Banach
completion with finite radius $\rho_*$. This is what
allows $\Phihol$ to be defined on
$\mathrm{Com}(H_k, V_k(\mathfrak{sl}_2))$ with explicit
Banach radius
$\rho_*(\mathrm{Com}) = |c| / (6|c_K(k)|)$
(Thm.~\ref*{opus:thm:PMC}(iii)).

\textbf{Maximal scope.} The inscribed scope
``non-logarithmic $C_2$-cofinite + VSKR+BGG-tempered
irrational coset'' is the \emph{union} of two strata:
\begin{itemize}
\item Stratum A: finite shadow-depth (class G/L/C/M in the
rational locus);
\item Stratum B: infinite-depth but Banach-tempered
(VSKR+BGG-tempered cosets).
\end{itemize}
Beyond this: logarithmic $W(p)$ is the nearest frontier
case (open); genuinely wild / irregular-singular algebras
(beyond the polar Stokes stratification) are outside the
programme's current reach. The maximal scope may be refined
upward by proving the $W(p)$ case (F1 frontier), but not
trivially.

\begin{hidden}[scope $\leftrightarrow$ Banach completion radius]
\label{hidden:cycle6}
The boundary between ``$\Phihol$ exists'' and ``$\Phihol$
does not exist'' is controlled by the Banach completion
radius $\rho_*(\cA) = |c|/\beta_\cA$: whenever $\rho_* > 0$,
there is a chain-level quasi-iso from the Banach completion
$\mathcal{B}_{\rho}(\cA)$ (for $\rho < \rho_*$) to
$\Zder(\cA)$, and $\Phihol(\cA)$ is defined as the bulk
HT-QFT whose bulk observables are this Banach completion.
When $\rho_* = 0$ (logarithmic $W(p)$ at $p \ge 2$, and
certain irregular-singular algebras), the bulk is not a
Banach space and the $\Etop{3}$-topological structure fails
structurally. The $\Phihol$-scope is exactly
$\{\cA : \rho_*(\cA) > 0\}$.
\end{hidden}
\end{heal}

%=======================================================================
\section{Construction of $\Phihol$ on $V_1(\mathfrak{sl}_2)$ from first
principles}
\label{sec:opus-construction}
%=======================================================================

The healed G3 identity is an \emph{existence-of-functor}
statement; this section constructs the functor at the
minimal non-trivial point. Input:
$(\cA, \omega, \mathrm{BL}) = (V_1(\mathfrak{sl}_2),
T_{\mathrm{Sug}}^{k=1}, \mathrm{Dirichlet})$.

\subsection{The boundary algebra}
$\cA = V_1(\mathfrak{sl}_2)$: the level-1 affine vertex
algebra of $\mathfrak{sl}_2$. Generators: three currents
$J^+, J^-, H$ with the level-1 affine OPEs
\begin{align}
  J^+(z) J^-(w) &\sim \frac{1}{(z-w)^2} + \frac{H(w)}{z-w}, \\
  H(z) J^\pm(w) &\sim \pm\frac{2 J^\pm(w)}{z-w}, \\
  H(z) H(w) &\sim \frac{2}{(z-w)^2}.
\end{align}
Non-critical level: $k+h^\vee = 1+2 = 3 \neq 0$. Central
charge from \eqref{eq:opus-ksug-c}:
$c_{\mathrm{Sug}}(1, \mathfrak{sl}_2) =
(1)(3)/(1+2) = 1$. So
$V_1(\mathfrak{sl}_2) = \mathrm{Vir}_{c=1}$-symmetric: the
Sugawara stress tensor has $c=1$. In fact
$V_1(\mathfrak{sl}_2) \simeq$ lattice VOA
$V_{\sqrt{2}\mathbb{Z}}$ (Frenkel--Kac).

\subsection{The conformal vector is inner at $k=1$}
Sugawara polynomial:
\begin{equation}
  T_{\mathrm{Sug}}(z) = \frac{1}{2 \cdot 3}
  \left({:}J^+ J^-{:} + {:}J^- J^+{:} +
  \tfrac{1}{2}{:}H H{:}\right)(z).
\end{equation}
BRST primitive in the Costello--Gaiotto bulk for
$\mathfrak{sl}_2$: introduce bulk antighosts
$\bar c^+, \bar c^-, \bar c_H$ dual to $J^+, J^-, H$. The
identity $J^a = [Q_{\mathrm{tot}}, \bar c^a]$ on
BRST cohomology (Costello--Gaiotto 2018) gives
\begin{equation}
  G^{(2)}(z) = \frac{1}{6}
  \left({:}\bar c^+ J^-{:} + {:}\bar c^- J^+{:} +
  \tfrac{1}{2}{:}\bar c_H H{:}\right)(z),
\end{equation}
verifying $T_{\mathrm{Sug}} = [Q_{\mathrm{tot}}, G^{(2)}]$
on cohomology by Leibniz.

\subsection{$\SCchtop$-heptagon action on the pair}
The pair $(\ChirHoch^\bullet(V_1(\mathfrak{sl}_2)),
V_1(\mathfrak{sl}_2))$ carries the $\SCchtop$-operad
action (\texttt{sc\_chtop\_heptagon.tex}): bulk braces
on boundary, bulk products from holomorphic $\Etop{2}$ on
$\mathrm{FM}_k(\mathbb{C})$, boundary products from
topological $\Etop{1}$ on $\mathbb{R}$. The closed-colour
generators are the chiral Deligne braces (HKR for
$V_1(\mathfrak{sl}_2) \simeq V_{\sqrt{2}\mathbb{Z}}$); the
open-colour generators are the boundary $E_1$-product on
$\mathbb{R}$ (Ayala--Francis); the mixed-colour operations
are the braces acting on $V_1(\mathfrak{sl}_2)$ by
$L_\infty$-brace formulas.

\subsection{Topologisation yielding $\Etop{3}$}
Tier~I: apply the Sugawara identity
$T_{\mathrm{Sug}} = [Q_{\mathrm{tot}}, G^{(2)}]$.
On BRST cohomology $H^\bullet_Q$ of the 3D BV complex,
$\partial_z$ is $Q_{\mathrm{tot}}$-exact:
$\partial_z \mathcal{O} =
[T_{\mathrm{Sug}}(0), \mathcal{O}] =
[[Q_{\mathrm{tot}}, G^{(2)}_{(0)}], \mathcal{O}] =
[Q_{\mathrm{tot}}, [G^{(2)}_{(0)}, \mathcal{O}]]$.
Therefore the factorisation algebra on $\mathbb{C}$
becomes locally constant on $H^\bullet_Q$, recognised
by Lurie HA Thm.~5.4.5.9 as an $\Etop{2}$. Dunn additivity
with the transverse $\Etop{1}$ from $\mathbb{R}$:
\begin{equation}
  H^\bullet_Q\bigl(\Zder(V_1(\mathfrak{sl}_2))\bigr)
  \;\in\; \mathrm{Alg}_{\,\Etop{3}}.
\end{equation}

\subsection{Chain-level original complex (class L)}
For $V_k(\fg)$ at $k \neq -h^\vee$, the class-L gauge
rectification of
\texttt{prop:sugawara-gauge-rectification} gives the
original-complex Tier~III statement. The explicit
$A_\infty$-coherence is computable at $V_1(\mathfrak{sl}_2)$:
the family $G = (G_1, G_2, G_3, \ldots)$ with
$G_1 = G^{(2)}_{(0)}$ (Sugawara antighost) and
$G_{\ge 2}$ given by higher-order Koszul-contraction
corrections supplying the coherence
$[m, G] = \partial_z$. Termination: at the level of
the shadow tower, class-L has $r_{\max} = 3$, so the
tower $G_n$ terminates at $n = 3$. This matches the
finite-depth argument
(\texttt{rem:topologization-scope}).

\subsection{Output: $\Phihol(V_1(\mathfrak{sl}_2))$}
The bulk HT-QFT is the Costello--Li abelian holomorphic
Chern--Simons on $\mathbb{C} \times \mathbb{R}$ for the
lattice $\sqrt{2}\mathbb{Z}$: the lattice VOA
$V_{\sqrt{2}\mathbb{Z}} \simeq V_1(\mathfrak{sl}_2)$ is the
boundary algebra, and the bulk factorisation algebra is
$\Zder(V_1(\mathfrak{sl}_2)) =
H^\bullet(V_{\sqrt{2}\mathbb{Z}}, V_{\sqrt{2}\mathbb{Z}})
\simeq$ lattice cohomology $H^\bullet(\mathrm{pt})[c]$,
equipped with $\Etop{3}$-structure from the double
Hochschild + transverse $\mathbb{R}$. This is the base
case, and it works chain-level on the original complex.

%=======================================================================
\section{Extension to $\cW_3$ and to $\cW_\infty[\mu]$}
\label{sec:opus-wstack}
%=======================================================================

\subsection{$\cW_3$ at generic $c$}
Input: $\cW_3^k$ at generic $c$, principal DS reduction of
$V_k(\widehat{\mathfrak{sl}}_3)$. Stress tower:
$\{T, W\}$ of spins $\{2, 3\}$. Antighost primitives
$G^{(2)}, G^{(3)}$ as in \S\ref{sec:opus-cycle-4}. Axiom T5
verified ($[G^{(2)}, G^{(3)}] = Q$-exact via lifting
Linshaw Thm.~7.1).

Tier~I output: $\Zder(\cW_3^k) \in \mathrm{Alg}_{\,\Etop{4}}$,
unconditional on cohomology at generic $c$.

Tier~II: transferred model unconditional; zero-differential
cohomology complex carries $\Etop{4}_\infty$-structure
unique up to $\infty$-quasi-iso.

Tier~III: weight-completed statement via DS-Hochschild
bridge to $V_k(\widehat{\mathfrak{sl}}_3)$; original-complex
statement OPEN (F3 frontier).

\subsection{$\cW_\infty[\mu]$ at generic $c, \mu$}
Input: $\cW_\infty[\mu]$, Linshaw's universal two-parameter
family. Infinite stress tower
$\{T^{(n)} : n \ge 2\}$. Axiom T5 at $N \ge 4$ POSTULATED.

Tier~I output: $\Zder(\cW_\infty[\mu]) \in
\mathrm{Alg}_{\,\Etop{\infty}}$ on cohomology, conditional
on T5 at all spins.

Tier~II: transferred model on cohomology complex,
conditional on T5.

Tier~III: OPEN.

\subsection{Convergence realisation}
At generic $c$, the three convergence criteria (a),(b),(c)
of Cycle~5 are equivalent and all hold; the
$\Etop{\infty}$-structure on
$\Zder(\cW_\infty[\mu])$ is realised chain-level via the
weight-filtered inverse limit, with each bar-weight window
$[w_{\min}, w_{\max}]$ computed in a finite truncation
$\cW_{w_{\max}-1}[\mu]$ by the finite-$N$ ladder.

\textbf{Minimal explicit test: spin-$\le 4$ truncation.}
$\cW_\infty[0]$ truncated to spin $\le 4$ carries the
stress tower $\{T, W_3, W_4\}$ (Zamolodchikov +
Kausch--Watts). Ladder output: $\Etop{5}$ (dimension
count: $2 + 2 + 1 = 5$). All three convergence criteria
verify by hand at this truncation
(\texttt{ex:w-infty-spin4-truncation}).

%=======================================================================
\section{Healed identifications}
\label{sec:opus-healed}
%=======================================================================

\begin{enumerate}
\item \textbf{Three-tier topologisation theorem} as in
Cycle~1, \texttt{heal:cycle1}: cohomological (I), transferred
(II), original-complex (III). Tiers I and II unconditional
across G/L/C/M (at weight-completed for class~M); Tier III
unconditional for class~L, frontier for class~M.

\item \textbf{Inner conformal vector is structural} as in
Cycle~2: innerness is a property of $(\cA, Q)$ requiring
existence of BRST primitive $G$ in the bulk BV complex. The
four standard families have inner conformal vectors at
non-critical level; logarithmic algebras fail by
$L_0$-nondiagonalisability.

\item \textbf{Critical-level failure is structural} as in
Cycle~3: at $k = -h^\vee$ the Sugawara denominator diverges,
the FF centre becomes commutative, and the
$\Etop{3}$-topological structure fails. Critical level
belongs to a different arithmetic programme (geometric
Langlands).

\item \textbf{$\cW_N$ ladder at $N=3$ via higher-spin
Casimir} as in Cycle~4: the $N=3$ case is the minimal
non-trivial ladder step beyond the base $N=2$; axiom T5 is
verifiable explicitly, the output is $\Etop{4}$, and
the $\cW_3$ algebra is the Zamolodchikov higher-spin
gravity boundary algebra.

\item \textbf{$\cW_\infty$ horizon is a genuine operadic
inverse limit} as in Cycle~5: each higher-spin Casimir is
algebraically independent, so the ladder does not
terminate; $\Etop{\infty}$ is defined as the inverse limit
of the stabilisation tower; three equivalent convergence
criteria (Gaberdiel--Gopakumar stabilisation,
$\infty$-categorical Dunn convergence, Yamada strong
convergence) realise it.

\item \textbf{Maximal scope of $\Phihol$ is
Banach-positive} as in Cycle~6: the scope is exactly
$\{\cA : \rho_*(\cA) > 0\}$, which includes non-logarithmic
$C_2$-cofinite + VSKR+BGG-tempered irrational cosets, and
excludes logarithmic $W(p)$ (F1 frontier) and
irregular-singular cases (beyond programme reach).

\item \textbf{Commit \texttt{a5640de}} is the $W(p)$
retraction, not an $H^2 = 0$ closure. The G3 programme
has no dependence on this commit for its topologisation
content.
\end{enumerate}

%=======================================================================
\section{Hidden structure}
\label{sec:opus-hidden}
%=======================================================================

\begin{hidden}[$\Phihol$ is Stage~2 of the two-stage
factorisation]
\label{hidden:phi-hol-stage-2}
The Vol~III two-stage factorisation
$\Phi_d = \mathrm{Sp}^{\mathrm{ch}}_{\Sigma_{d-1}, C} \circ
\Phi^{\mathrm{FA}}_d$ gives Stage~1 a canonical
$E_d$-holomorphic factorisation algebra on a CY-$d$ variety,
and Stage~2 specialises via factorisation-homology
pushforward $\int_{\Sigma_{d-1}}$ to a reference curve $C$.
The output of Stage~2 is an $E_1$-chiral algebra on $C$
equipped with the 3D HT-QFT bulk
$\Zder(\cdot)$. The functor $\Phihol$ is Stage~2 of this
factorisation: $\Phihol = \mathrm{bulk} \circ
\mathrm{Sp}^{\mathrm{ch}}_{\Sigma_{d-1}, C}$. The two
colours of $\SCchtop$ are the two stages: closed colour =
Stage~1 $E_d$-holomorphic, open colour = Stage~2 $E_1$-chiral
(after specialisation). This observation is in
\texttt{sc\_chtop\_heptagon.tex}
Rem.~\texttt{rem:heptagon-two-stage-CY-to-chiral}.
\end{hidden}

\begin{hidden}[Virasoro climax is $\Phihol$ at $\cA = \Vir_c$]
\label{hidden:virasoro-climax}
The 3D quantum gravity cluster of Vol~II is not a
separate mathematical structure from $\Phihol$; it is
$\Phihol(\Vir_c)$ for $c \neq 0$ non-critical. The
Brown--Henneaux identification $c = 3\ell/(2G_N)$ converts
the abstract central charge into a geometric input (AdS
radius $\ell$, Newton constant $G_N$); the Virasoro boundary
algebra determines the partition function and BTZ spectrum;
the bulk is read off as $\Zder(\Vir_c)$. This is the
content of Corollary~P2.1 (memory file
\texttt{platonic\_ideal\_reconstituted\_2026\_04\_17.md}).
The identification ``bulk 3D QG partition function = Vol~I
shadow-tower generating function'' is a computation
mediated by $\Phihol$, not a shadow.
\end{hidden}

\begin{hidden}[Topologisation tower $=$ higher-spin gravity
dual]
\label{hidden:topo-tower-higher-spin}
The ladder
$\Etop{3} \to \Etop{4} \to \cdots \to \Etop{\infty}$
produced by iterated Sugawara is the bulk dual of the
Gaberdiel--Gopakumar higher-spin AdS$_3$/CFT$_2$
correspondence. The CFT$_2$ boundary is $\cW_N$ or
$\cW_\infty[\mu]$; the AdS$_3$ bulk is Vasiliev higher-spin
gravity of spin range $[2, N]$ or $[2, \infty]$. The
$\Etop{N+1}$-operadic upgrade is the algebraic shadow of
the $(N+1)$-dimensional transverse geometric data of the
Vasiliev theory.
\end{hidden}

\begin{hidden}[Dunn additivity on the topologised direction
only]
\label{hidden:dunn-topologised}
The Dunn-additivity step
$\Etop{p} \otimes \Etop{q} \simeq \Etop{p+q}$ applies only
when both factors are topological. The chiral direction is
not originally topological; it is $\Ech{2}$-chiral by
chiral Deligne. Sugawara is what converts it to
$\Etop{2}$. Without this conversion, the naive Dunn
$\Ech{2} \otimes \Etop{1} \ne \Etop{3}$ fails (AP,
\texttt{ap:topologization-healed}). The antipattern
precisely says: one must topologise \emph{before} applying
Dunn. The Sugawara mechanism is the topologisation step; it
is not a parallel structure to Dunn, it is what unlocks Dunn.
\end{hidden}

%=======================================================================
\section{Open frontiers}
\label{sec:opus-open}
%=======================================================================

\begin{openfrontier}[F3: class-M chain-level on the original
complex]
\label{open:F3}
Prove that $\Zder(\cW_N^k)$ (and more generally
$\Zder(\cA)$ for class-M algebras) carries a chain-level
$\Etop{N+1}$-structure on the \emph{original} (uncompleted,
direct-sum) complex, not merely on the weight-completed
one. Obstruction: the iterated-Sugawara
$A_\infty$-coherence tower need not terminate; the
class-L finite-shadow-depth argument does not apply.
Current scope: weight-completed holds
(\texttt{prop:bv-bar-class-m-weight-completed}); original
complex is F3.

Strategic routes: (a) find an explicit combinatorial
formula for the $A_\infty$-tower and prove it terminates
at $r_{\max}(\cW_N)$; (b) prove that weight-completion is
essentially surjective in the derived category, which
would transport the weight-completed result to the
original complex up to equivalence; (c) find a direct 3D
BV BRST construction of
$\Phihol(\cW_N)$ independent of DS-reduction and show that
its chain-level $\Etop{3}$-structure restricts to the
original complex directly.
\end{openfrontier}

\begin{openfrontier}[T5 at high spin]
\label{open:T5}
Prove $[G^{(n)}, G^{(m)}] = Q_{\mathrm{tot}}$-exact for all
$n, m \ge 4$ from the Costello--Gaiotto bulk BV-BRST
structure on the 3D holomorphic-topological quantisation of
a chiral algebra with higher-spin stress tower, not by
appeal to the classical $\cW_\infty[\mu]$ Poisson algebra.
Until this is established, the $\Einfty$-specialisation
Theorem $\cW_\infty \to \Etop{\infty}$ is conditional.
\end{openfrontier}

\begin{openfrontier}[F1: logarithmic $W(p)$ tempering]
\label{open:F1}
Prove (or disprove at each $p$) that the logarithmic
triplet $W(p)$ admits a tempered-stratum membership in the
$\Phihol$-scope. The Zhu-bounded-Massey route fails by
commit \texttt{a5640de} (Gurarie 1993, Flohr 1996); a new
mechanism is needed. Candidates: (a) direct character-
amplitude bound via Adamovi\'c--Milas fusion rules; (b)
logarithmic 3D gravity bulk construction and explicit
BRST primitive for $\omega = T_{W(p)}$ in its log-BRST
complex.
\end{openfrontier}

\begin{openfrontier}[Class-L original-complex TIER III:
$A_\infty$-tower termination proof]
\label{open:classL-tierIII}
The class-L Tier III statement is inscribed as proved via
\texttt{prop:sugawara-gauge-rectification}, but the
explicit $A_\infty$-coherence tower requires finite shadow
depth ($r_{\max} = 3$ for class~L) to terminate. A purely
chain-level proof of the termination (without appeal to
the shadow-depth-finiteness abstract argument) is not yet
inscribed. Strengthening the Tier III statement to a
direct chain-level construction would improve the
epistemic standing from
``proved via finite-depth argument'' to ``proved via
explicit chain-level formula''.
\end{openfrontier}

%=======================================================================
\section{Cross-volume harmonies}
\label{sec:opus-harmonies}
%=======================================================================

\subsection{Vol I harmonies}
\begin{itemize}
\item The G3 topologisation base is inscribed in Vol I
(\texttt{en\_koszul\_duality.tex} \S\ref*{opus:sec:V1topo}),
proving
Thm.~\ref*{opus:thm:topologization} for $V_k(\fg)$,
$k \neq -h^\vee$, with explicit three-tier layer structure
(Tier~I, II, III).
\item The iterated Sugawara
$\Einfty$-tower is inscribed in Vol I
Thm.~\ref*{opus:thm:E-inf-top}, matching the Vol II
\texttt{e\_infinity\_topologization.tex} ladder. The two
volumes agree on scope: Tier~I unconditional; Tier~II
unconditional; Tier~III class-L unconditional,
class-M weight-completed, class-M original-complex open.
\item The G3 functor $\Phihol$ is the Vol~II operationalisation
of Vol~I's Theorem~A (bar-cobar backbone): the
bulk-boundary adjunction $\mathrm{bulk}: \cA \mapsto
\Zder(\cA)$ and $\mathrm{bdry}: T_\cA \mapsto
\Obs^\partial(T_\cA)$ is the chiral instantiation of
$\Omega^{\mathrm{ch}} \dashv B^{\mathrm{ch}}$ at the
factorisation-algebra level (HT-QFT ambient).
\end{itemize}

\subsection{Vol II harmonies}
\begin{itemize}
\item $\Phihol$ is the master theorem
(\texttt{universal\_holography\_functor.tex}
Thm.~\ref*{opus:thm:UHF}), with explicit scope
``non-logarithmic $C_2$-cofinite + VSKR+BGG-tempered''
(Cycle~6).
\item The $\Einfty$-ladder
(\texttt{e\_infinity\_topologization.tex}) specialises to
Virasoro ($\Etop{3}$), $\cW_N$ ($\Etop{N+1}$),
$\cW_\infty$ ($\Etop{\infty}$). Scope is precise: $N \le 3$
unconditional (axiom T5 vacuous or reducible to one
bracket), $N \ge 4$ conditional on T5.
\item The class-M weight-completed status comes from the
DS-Hochschild bridge
(\texttt{chd-ds-hochschild}), which reduces class-M to
class-L under Drinfeld--Sokolov.
\end{itemize}

\subsection{Vol III harmonies}
\begin{itemize}
\item The $\SCchtop$ two-colour structure is the
algebraic encoding of the Vol~III two-stage factorisation
$\Phi_d = \mathrm{Sp}^{\mathrm{ch}}_{\Sigma_{d-1}, C} \circ
\Phi^{\mathrm{FA}}_d$: closed colour = Stage~1
$E_d$-holomorphic FA, open colour = Stage~2
$E_1$-chiral after specialisation.
\item The B-row ($\kappa + \kappa^! = 8$) Mukai-enhanced K3
Heisenberg chiral algebra has a \emph{non-inner} conformal
vector: it enters $\Phihol$ via a different route (direct
Borcherds-Heegner reciprocity) rather than Sugawara
topologisation.
\item The Universal Trace Identity
(memory file \S4) closes across the three volumes: Vol I
ghost-scope $\mathrm{tr}_{\mathrm{ghost}}(Q^2)$, Vol II
Pentagon-scope $\mathrm{tr}_{\mathrm{Pentagon}}$, Vol III
Borcherds-scope $\omega_{\mathrm{Borcherds}}$ all equal
$c_N(0)/2$ at each $N \in \{1,2,3,4,6\}$ and the
half/quarter-integer continuations.
\end{itemize}

%=======================================================================
\section{Cycle log: sequence of attack-heal turns}
\label{sec:opus-log}
%=======================================================================

\textbf{Cycle 1 (Attacks a, g).}
Attack: identification-vs-lift + commit \texttt{a5640de}
question. Healed by three-tier layer structure + direct
\texttt{git show} verification that \texttt{a5640de} is the
$W(p)$ retraction. Hidden: three tiers are different
theorems about different objects.

\textbf{Cycle 2 (Attack b).}
Attack: is ``inner'' structural or decorative? Healed:
inner requires BRST primitive $G$ in the bulk BV complex;
the identity $T = [Q, G]$ in cohomology is what lets
holomorphic translations become $Q$-exact. Non-inner
algebras (B-row / Borcherds) enter via a different route.

\textbf{Cycle 3 (Attack c).}
Attack: regularity + critical-level. Healed: chain-level
single-contour normally ordered for $k \neq -h^\vee$,
structural failure at critical level (FF centre becomes
commutative). Hidden: critical vs non-critical is the
Langlands wall.

\textbf{Cycle 4 (Attack d).}
Attack: $\cW_N$ chain-level + axiom T5 scope. Healed:
$N = 3$ is the minimal non-trivial test; T5 reduces to one
bracket $[G^{(2)}, G^{(3)}]$, verified via Leibniz +
Linshaw Thm~7.1; $N \ge 4$ requires T5 postulated, with
the Linshaw classical-Poisson identity as evidence. Output:
$\Etop{4}$ for $\cW_3$.

\textbf{Cycle 5 (Attack e).}
Attack: $\cW_\infty$ horizon -- genuine or finite
termination? Healed: each higher-spin Casimir is
algebraically independent (Linshaw commutant theorem), so
the tower does not stabilise. $\Etop{\infty}$ is the
inverse limit in $\mathrm{Op}_\infty^{\otimes}$; three
equivalent convergence criteria realise it chain-level at
generic $c$.

\textbf{Cycle 6 (Attack f).}
Attack: scope of unconditional $\Phihol$. Healed: maximal
scope is $\{\cA : \rho_*(\cA) > 0\}$ = Banach-positive
algebras, union of $C_2$-cofinite strata and VSKR+BGG
tempered irrational cosets. $W(p)$ (logarithmic) and
irregular-singular cases are outside.

\textbf{Summary of healed $\Phihol$ status.}
\begin{center}
\begin{tabular}{lccc}
\toprule
Class & Tier~I (coh) & Tier~II (transf) & Tier~III (orig) \\
\midrule
G (Heis/lattice) & Unconditional & Unconditional & Unconditional \\
L ($V_k(\fg)$, $k\neq-h^\vee$) & Unconditional & Unconditional & Unconditional (class~L) \\
C ($\beta\gamma$, $bc$) & Unconditional & Unconditional & Unconditional (boson.\ G) \\
M-wc (Vir, $\cW_N$, wt-comp.) & Unconditional & Unconditional & Unconditional \\
M-oc (Vir, $\cW_N$, orig.~cplx) & Unconditional & Unconditional & OPEN (F3) \\
B-row (Mukai K3 Heis) & Unconditional & Unconditional & Unconditional (direct) \\
$W(p)$ (logarithmic) & OPEN (F1) & OPEN & OPEN \\
\bottomrule
\end{tabular}
\end{center}

\textbf{Summary of $\Einfty$-ladder status.}
\begin{center}
\begin{tabular}{lcc}
\toprule
Depth $N$ & Ladder output & Status \\
\midrule
2 (Virasoro, affine KM) & $\Etop{3}$ & Unconditional (axiom T5 vacuous) \\
3 ($\cW_3$) & $\Etop{4}$ & Unconditional (T5 at $n,m \in \{2,3\}$ verifiable) \\
$N \ge 4$ ($\cW_N$) & $\Etop{N+1}$ & Conditional on T5 at $n, m \ge 4$ \\
$\infty$ ($\cW_\infty[\mu]$) & $\Etop{\infty}$ & Conditional on T5 all spins \\
\bottomrule
\end{tabular}
\end{center}

\end{document}
